\noindent{\Large\bf Introduction}

These notes are the very late written version of a series of lectures
given at the Trieste Summer School in 1995.  The aim was to provide an
elementary introduction to the work of Seiberg and Witten on exact
results concerning $N=2$ supersymmetric extensions of Quantum
Chromodynamics.  We wanted to provide, in a single place, all the
background material necessary to study their work in detail.  We had
in mind graduate students who have already gone through their Quantum
Field Theory course, but we do not expect much more background to
follow these lectures. We have done our best to make the treatment
pedagogical.  In some sections, we have heavily drawn on previous
reviews, for instance in the treatment of supersymmetry we have
followed Bagger and Wess \cite{WB}, and in the theory of monopoles we
have used the reviews by Goddard and Olive \cite{GO} and by S. Coleman
\cite{coleman1,coleman2}.  These sources provide excellent
presentations of these topics, and we had no compelling reason to try
to make improvements on their presentation.  It is also quite obvious
that we have drawn heavily on the original papers of Seiberg and
Witten \cite{SW-I,SW2}, but we have tried to provide the necessary
tools to make their reading more accessible to interested students
and/or researchers. Recently there have been other lecture notes
published at a more advanced level, where one can find more details
and also the connection with String Theory (see, for instance, the
notes by W. Lerche \cite{Lerche}).  We would also like to mention that
we have not tried to be exhaustive in quoting all the literature on
the subject.  A more complete reference list can be found in
\cite{Lerche}.  The reference list is intented to provide a guidance
to enter the vastly growing literature on duality in String Theory and
Field Theory.  We apologize to those authors who may be offended by
not finding their works referenced.

These notes are divided into four sections, with each section further
subdivided into several subsections. Section 1 is devoted to an
introduction to monopoles in gauge theories. We start with a
discussion of the Dirac monopole and the idea of charge quantization,
and then describe the 't Hooft-Polyakov monopole in gauge theories
with spontaneous symmetry breaking. Then we introduce the notion of
Bogomol'nyi bound and the BPS states. After this, we describe the
topological classification of monopoles and then describe the
Montonen-Olive conjecture of electric-magnetic duality. We end this
section with a description of how, in the presence of a $\theta$-term
in the Lagrangian, the electric charge of a monopole is shifted by its
magnetic charge. Section 2 is devoted to an introduction to
supersymmetric gauge theories. First we describe the supersymmetry
algebra and its representations without and with central charges and
discuss its local realizations in terms of superfields. Then we
construct $N=1$ Lagrangians and finally, $N=2$ supersymmetric
Lagrangians with both gauge multiplets and matter multiplets
(hypermultiplets). At the end, we calculate the $N=2$ central charge
both in the pure gauge theory, as well as in the theory with matter
and establish its relation to the BPS bound. Having built the
foundations in the first two sections, in section 3 we describe the
Seiberg-Witten analysis of the $N=2$ pure gauge theory with gauge
group $SU(2)$. In the first two subsections, we discuss the
parametrization of the moduli space and the breaking of R-symmetries.
Then we describe how the chiral $U(1)$ anomaly of the theory can be
used to obtain the one-loop form of the low-energy effective action.
The rest of the section is devoted to finding the exact low-energy
effective action by using duality and the singularity structure on the
moduli space of the theory. We express the exact solution in terms of
complete elliptic integrals. In section 4, we briefly describe the
Seiberg-Witten analysis of $N=2$ $SU(2)$ gauge theory with $N_f$
matter fields. After a discussion of the general features of these
theories, we describe how the duality group is no longer pure
$SL(2,Z)$. Then we describe the singularity structure on the moduli
space of these theories and sketch the procedure for obtaining the
exact solutions. Our aim in section 4 is to give a flavour of the
analysis of theories with matter and, for a deeper understanding, the
interested reader is referred to the original work of Seiberg and
Witten.

\section{Magnetic Monopoles in Gauge Theories}

In this section, we begin by reviewing the proporties of the Dirac
monopole and the idea of charge quantization. Then we describe the
magnetic monopoles and dyons which arise in non-Abelian gauge theories
with spontaneous symmetry breaking and discuss their general
properties. We also introduce the notion of the Bogomol'nyi bound and
BPS states. In the last two subsections, we describe the
Montonen-Olive conjecture of electric-magnetic duality and Witten's 
argument about how the presence of the $\theta$-term in the Lagrangian
modifies the monopole and dyon electric charges. For a more detailed
discussion of most of the material in this section, the reader is
referred to the review article by Goddard and Olive \cite{GO}, and to
the lecture notes by Coleman \cite{coleman1,coleman2}.

\subsection{Conventions and Preliminaries}

Let us start by stating our conventions: we always take $c=1$, and
almost always $\hbar=1$, except when it is important to make a
distinction between classical and quantum effects. For index
manipulations, we use the flat Minkowsky metric $\eta$ of signature
$\{ +,-,-,-\}$.  Moreover, we choose units in which Maxwell's
equations take the form:
\begin{equation}
\begin{array}{l}
\vec{\nabla}\cdot\vec{E}=\rho\,,\qquad\vec{\nabla}\times\vec{B}-
\partial \vec{E}/\partial t =\vec{j}\,,\\
\vec{\nabla}\cdot\vec{B}=0\,,\qquad \vec{\nabla}\times\vec{E} +
\partial\vec{B}/\partial t = 0\,.
\label{I}
\end{array}
\end{equation}
In these units, a factor of $(4 \pi)^{-1}$ appears in Coulomb's law:
For a static point-like charge $q$ at the origin, we have $\rho = q 
\delta^{3}(\vec{r})$. Integrating the first Maxwell equation over a
sphere of radius $r$ and using spherical symmetry, we get
$\int_{S^2} \vec{E} \cdot d \vec{s}=4\pi r^2 E(r)=q$. Hence,
$\vec{E}=q\vec{r}/4\pi r^{3}$, as stated. The electrostatic
potential $\phi$ defined by $\vec{E}=-\vec{\nabla}\phi$ is given by 
$\phi = q / 4 \pi r$. 

In relativistic notation, one introduces the four-potential 
$A^{\mu}= \{\phi, \vec{A}\}$. The electric and magnetic fields are
defined as components of the corresponding field strength tensor
$F_{\mu\nu}$ as follows:
\begin{eqnarray}
F_{\mu \nu} & = &\partial_{\mu}A_{\nu} - \partial_{\nu}A_{\mu}\,,
\nonumber \\
F_{0i} & = &\partial_{0}A_{i} - \partial_{i}A_{0} =
- \partial_{0}A^{i} - \partial_{i}A^{0} = E^{i}\,,\nonumber \\
F_{ij} & = & \partial_{i}A_{j} - \partial_{j}A_{i} =
- (\partial_{i}A^{j} - \partial^{j}A_{i}) =
-\epsilon_{ijk} B^{k}\,,
\label{III}
\end{eqnarray}
so that $B^i=-\frac{1}{2}\epsilon^{ijk}F_{jk}=-\frac{1}{2}
\epsilon^{oi\mu\nu}F_{\mu\nu}$. The dual field strength tensor is
given by   
$$
\widetilde F^{\mu \nu} = \frac{1}{2} \varepsilon^{\mu \nu \alpha
\beta} F_{\alpha \beta}\,,
$$
with $\varepsilon^{0123}=+1$. In component notation, we have 
\begin{equation}
F^{\mu \nu} =
\left(
\begin{array}{cccc}
0 & -E_{x} & -E_{y} & -E_{z}  \\
E_{x} & 0 & -B_{z} & B_{y} \\
E_{y} & B_{z} & 0 & -B_{x} \\
E_{z} & -B_{x} & B_{y} & 0
\end{array}
\right)\,,
\qquad
\widetilde F^{\mu \nu} =
\left(
\begin{array}{cccc}
0 & -B_{x} & -B_{y} & -B_{z} \\
B_{x} & 0 & E_{z} & -E_{y} \\
B_{y} & -E_{z} & 0 & E_{x} \\
B_{z} & E_{y} & -E_{x} & 0
\end{array}
\right)\,.
\label{IV}
\end{equation}

In terms of the electric four-current $j^{\mu}=\{\rho,\vec{j}\}$, the
Maxwell's equations take the compact form
\begin{equation}
\partial_\nu F^{\mu\nu}=-j^\mu\,,\qquad\partial_\nu\,\widetilde
F^{\mu\nu}=0\,. \label{V}
\end{equation}
Note that when $j^\mu =0$, the above equations are invariant under the
replacement $\vec{E} \rightarrow \vec{B}, \vec{B} \rightarrow
-\vec{E}$. This is referred to as the electric-magnetic duality
transformation. In the presence of electric sources, this
transformation is no longer a symmetry of Maxwell's equations. In
order to restore the duality invariance of these equations for
non-zero $j^{\mu}$, Dirac \cite{Dirac} introduced the magnetic
four-current $k^{\mu}=\{\sigma ,\vec{k}\}$ and modified Maxwell's
equations to
\begin{equation}
\partial_\nu F^{\mu\nu}=-j^\mu\,,\qquad\partial_\nu\,
\widetilde F^{\mu\nu}= -k^\mu \,.   
\label{VI}
\end{equation}
The above equations are now invariant under a combined duality
transformation of the fields and the currents which can be written as 
\begin{equation}
F\rightarrow \widetilde F\,,\,\widetilde F\rightarrow-F\,;\qquad 
j^\mu\rightarrow k^\mu\,,\, k^\mu\rightarrow -j^\mu\,.
\label{emduality}
\end{equation}
Note that the full invariance group of the equations (\ref{VI}) is
larger than this discrete duality. In fact, they are invariant under a
continuous $SO(2)$ group which rotates the electric and magnetic
quantities into each other. 

For point-like electric and magnetic sources, the current densities 
can be written as 
\begin{eqnarray}
j^{\mu}(x) & = & \sum_{a} q_{a} \, \int \, dx_{a}^{\mu}
\delta^{4} (x - x_{a})\,, \nonumber \\
k^{\mu}(x) & = & \sum_{a} g_{a} \, \int \, dx_{a}^{\mu}
\delta^{4} (x - x_{a})\,.\nonumber
\end{eqnarray}
A particle of electric charge $q$ and magnetic charge $g$ experiences
a Lorentz force given by
$$
m\,\frac{d^2x^\mu}{d\tau^2}=(q F^{\mu\nu}+g\, \widetilde F^{\mu\nu})
\frac{dx_\nu}{d\tau}\,.
$$

Although, at the level of quations of motion, Dirac's modification of
Maxwell's theory seems trivial, it has highly non-trivial consequences
in quantum theory. One way of realizing the problem is to note that
the vector potential $\vec{A}$ is indispensable in the quantum
formulation of the theory. On the other hand, equations
$\vec{B}=\vec{\nabla}\times \vec{A}$ and $\vec{\nabla}\cdot\vec{B}\ne
0$ are not compatible unless the vector potential $\vec{A}$ has
singularities and, hence, is not globally well-defined. It turns out
that these singularities are gauge dependent and, therefore, their
presence should not be experimentally  detectable. In the classical
theory, which can be formulated in terms of $\vec B$ alone, this
requirement is trivially satisfied. However, in quantum theory, it
leads to the important phenomenon of the quantization of electric
charge. In the following, we will first give a semiclassical
derivation of this effect and then describe a derivation based on the
notion of the Dirac string. 
 
\subsection{A Semiclassical Derivation of Charge Quantization}

Consider a non-relativistic charge $q$ in the vicinity of a magnetic
monopole of strength $g$, situated at the origin. The charge $q$
experiences a force $m\ddot{\vec{r}}=q\dot{\vec{r}}\times\vec{B}$,
where $\vec{B}$ is the monopole field given by $\vec{B}= 
g\vec{r}/4\pi r^3$. The change in the orbital angular momentum of the
electric charge under the effect of this force is given by 
$$
\frac{d}{dt}\left(\,m\vec{r}\times\dot{\vec{r}}\right) =
m\vec{r}\times\ddot{\vec{r}}= \frac{qg}{4 \pi r^{3}}\,
\vec{r}\times\left(\dot{\vec{r}}\times\vec{r}\right)= 
\frac{d}{dt}\left(\frac{qg}{4\pi}\,\frac{\vec{r}}{r}\right)\,.
$$
Hence, the total conserved angular momentum of the system is
\begin{equation}
\vec{J}=\vec{r}\times m\dot{\vec{r}}-\frac{qg}{4\pi}\,
\frac{\vec{r}}{r}\,.
\label{XI}
\end{equation}
The second term on the right hand side (henceforth denoted by
$\vec{J}_{em}$) is the contribution coming from the elecromagnetic
field.  This term can be directly computed by using the fact that the
momentum density of an electromagnetic field is given by its Poynting
vector, $\vec{E}\times\vec{B}$, and hence its contribution to the
angular momentum is given by
$$
\vec{J}_{em} = \int \, d^{3}x \, \vec{r} \times
(\vec{E} \times \vec{B}) = \frac{g}{4 \pi} \,
\int \, d^{3} x \, \vec{r} \times
\left( \vec{E} \times \frac{\vec{r}}{r^{3}} \right)\,.
$$
In components,
\begin{eqnarray}
J^i_{em} & = & \frac{g}{4 \pi} \, \int \, d^{3}x E^{j}
\left( \delta_{ij} - \frac{x_i x_j}{r^{2}} \right )
\frac{1}{r} = \frac{g}{4 \pi} \, \int \, d^{3}x E^{j}
\partial_{j} (\hat{x}^{i}) \nonumber \\
& = & \frac{g}{4 \pi} \int_{S^2} \hat{x}^i \vec{E}\cdot \vec{ds} 
- \frac{g}{4 \pi} \, \int \, d^{3}x \vec{\nabla} \cdot
\vec{E} \hat{x}^{i}\,. 
\label{XIV}
\end{eqnarray}
When the separation between the electric and magnetic charges is
negligible compared to their distance from the boundary $S^2$, the
contribution of the first integral to $\vec{J}_{em}$ vanishes by
spherical symmetry. We are therefore left with   
\begin{equation}
\vec{J}_{em}= -\frac{gq}{4 \pi} \hat{r}
\label{Jem}\,.
\end{equation}

Returning to equation (\ref{XI}), if we assume that orbital angular
momentum is quantized. Then it follows that 
\begin{equation}
\frac{qg}{4 \pi} = \frac{1}{2} n \hbar\,,
\label{XII}
\end{equation}
where $n$ is an integer. Note that in the above we have assumed the
total angular momentum of the charge-monopole system to be quantized
in half-integral units. This is a strange assumption considering that
we did not have to treat the electrically charged particle or the
monopole as fermios.  both of the components are bosonic. However, it
turns out that this actually is the case and that the situation does
not contradict the spin-statistics theorem
\cite{Goldhaber,Wilczek}. We will not discuss this issue further but
remark that the derivation of the same equation presented in the next
subsection does not depend on this assumption.

Equation (\ref{XII}) is the Dirac charge quantization
condition. It implies that if there exists a magnetic monopole of
charge $g$ somewhere in the universe, then all electric charges are
quantized in units of $2\pi\hbar/g$. If we have a number of purely
electric charges $q_i$ and purely magnetic charges $g_j$, then any
pair of them will satisfy a quantization condition:
\begin{equation}
\frac{q_{i}g_{j}}{4 \pi \hbar} = \frac{1}{2} n_{ij}
\label{XIII}
\end{equation}
Thus, any electric charge is an integral multiple of $2\pi\hbar/g_j$.
For a given $g_j$, let these charges have $n_{0j}$ as the highest
common factor. Then, all the electric charges are multiples
of $q_{0}=n_{0j}2\pi\hbar /g_{j}$. Note that $q_0$ itself may not
exist in the spectrum. Similar considerations apply to the
quantization of magnetic charge. 

Till now, we have only dealt with particles that carry either an
electric or a magnetic charge. Let us now consider dyons, {\it i.e.,}
particles that carry both electric and magnetic charges. Consider two
dyons of charges $(q_1,g_1)$ and $(q_2,g_2)$. For this system, we can
repeat the calculation of $\vec{J}_{em}$ by following the steps in
(\ref{XIV}) where now the electromagnetic fields are split as
$\vec{E}=\vec{E}_1+\vec{E}_2$ and $\vec{B}=\vec{B}_1+ \vec{B}_2$.  The
answer is easily found to be

\begin{equation}
\vec{J}_{em}=-\frac{1}{4\pi}\left(q_1 g_2 - q_2 g_1\right)\hat{r} 
\label{Jemd}
\end{equation}
The charge quantization condition is thus generalized to 
\begin{equation}
\frac{q_1 g_2 - q_2 g_1}{4\pi\hbar}=\frac{1}{2} n_{12}
\label{DSZ}
\end{equation}
This is referred to as the Dirac-Schwinger-Zwanziger condition
\cite{Schwinger,Zwanziger}. This condition is invariant under the
$SO(2)$ transformation $(q+ig)\rightarrow e^{i\phi}(q+ig)$ which is
also a symmetry of (\ref{VI}). 

\subsection{The Dirac String}

In the following, we present a more rigorous derivation of the Dirac
quantization condition which is based on the notion of a Dirac
string. Let $\vec{B}_{mon}$ denote the magnetic field around a
monopole. Since $\vec{\nabla}\cdot\vec{B}_{mon}\ne 0$, it is not
possible to construct a well-defined $\vec{A}_{mon}$ such that
$\vec{B}_{mon}=\vec{\nabla} \vec{A}_{mon}$. To overcome this problem,
Dirac intruduced a  semi-infinite solenoid (or string) running from
$(0,0,-\infty)$ to the monopole position, $(0,0,0)$. This solenoid
carries a magnetic field $\vec{B}_{sol}=g\theta (-z)\delta(x)
\delta(y)\hat{z}$ which also is not divergence free. However, the
total magnetic field,   
\begin{equation}
\vec{B}=\vec{B}_{mon}+\vec{B}_{sol}=\frac{g}{4\pi r^{2}}+g\theta(-z)
\delta(x)\delta(y)\hat{z}\,,
\end{equation}
satisfies $\vec{\nabla}\cdot\vec{B}=g\delta(\vec{r})-g\delta(\vec{r}) 
=0$. It is therefore possible to construct a non-singular $\vec{A}=
\vec{A}_{mon}+\vec{A}_{sol}$ corresponding to the monopole-solenoid
system. In fact, the singular $\vec{A}_{sol}$ associated with the
Dirac string is used to cancel the singularity in $\vec{A}_{mon}$. The
position of the singularity in $\vec{A}_{mon}$, and therefore the
position of the Dirac string, can be shifted by singular gauge
transformations. Since the Dirac string is an artificial construct, it
should be unobservable and should not contribute to any physical
process. However, in an Aharonov-Bohm experiment, the presence of the
string can affect the wavefunction of an electric charge by
contributing to its phase. This resulting Aharonov-Bohm phase along a
contour $\Gamma$ encircling the string and enclosing an area $S$ can
be easily computed to be
$$
e \oint_{\Gamma} \vec{A}_{sol} \cdot d \vec{l} = e \int_S
\vec{B}_{sol} \cdot d\vec{s} = e g\,.
$$
Here, $e=q/\hbar$ is the electromagnetic coupling constant. This
phase, and therefore the Dirac string, is unobservable provided 
$$
e g = 2\pi n\,,
$$
which is again the Dirac charge quantization condition.

% ==== FIGURE_SLOT_1 ====
\subsection{The Georgi-Glashow Model: A Simple Theory with Monopoles}

Till now we have been working in the framework of particle mechanics
where both electric and magnetic charges are point-like objects and
are introduced by hand into the theory. However, in field theory,
these objects can also arise as solitons which are non-trivial
solutions of the field equations with localized energy density. If the
gauge field configuration associated with a soliton solution has a
magnetic character, the soliton can be identified as a magnetic
monopole. The simplest solitons are found in the Sine-Gordon theory,
which is a scalar field theory in $1+1$ dimensions. However, by
Derrick's theorem, a scalar field theory in more than two dimensions
cannot support static finite energy solutions. This is basically due
to the fact that because of the non-trivial structure of the soliton
at large distances, the total energy of the configuration
diverges. This situation can be cured by the addition of gauge fields
to the theory. Thus, a scalar theory with gauge interactions in four 
dimensions can admit static finite energy field configurations. The
stability of such solitonic configurations are often related to the
fact that they are characterized by conserved topological charges. 
For a more detailed discussion of these issues, the reader is referred
to \cite{GO, coleman1, coleman2}. In the following, we describe the
Georgi-Glashow model which is a simple theory in $3+1$ dimensions with
monopole solutions.  

The Georgi-Glashow model is a Yang-Mills-Higgs system which contains 
a Higgs multiplet $\phi^a\,(a=1,2,3)$ transforming as a vector in the
adjoint representation of the gauge group $SO(3)$, and the gauge
fields $W_\mu=W^a_\mu T^a$. Here, $T^a$ are the hermitian generators
of $SO(3)$ satifying $[T^a,T^b]=if^{abc}T^c$. In the adjoint
representaion, we have $(T^a)_{bc}=-if^a_{bc}$ and, for $SO(3)$,
$f^{abc}=\epsilon^{abc}$. The field strength of $W_\mu$ and the
cavariant derivative on $\phi^a$ are defined by 
\begin{eqnarray}
G_{\mu \nu} &=& \partial_{\mu}W_{\nu} - \partial_{\nu} W_{\mu} +
ie [W_{\mu},W_{\nu}]\,,
\nonumber\\
D_\mu\phi^a &=& \partial_\mu\phi^a - e\epsilon^{abc} W^b_\mu \phi^c\,.
\end{eqnarray}
The minimal Lagrangian is then given by
\begin{equation}
{\cal L}  =  - \frac{1}{4} G^{a}_{\mu \nu} G^{a \mu \nu} +
\frac{1}{2} D^{\mu}\phi^a D_{\mu}\phi^a - V(\phi)\,,
\label{GGmodel}
\end{equation}
where,
\begin{equation}
V(\phi)=\frac{\lambda}{4}\left(\phi^2 - a^2 \right)^2\,.
\label{GGpot}
\end{equation}
The equations of motion following from this Lagrangian are 
\begin{equation}
(D_\nu G^{\mu \nu})^a=-e\,\epsilon^{abc}\,\phi^b\,(D^\mu\phi)^c,
\qquad D^\mu D_\mu \phi^a =-\lambda\phi^a (\phi^{2} - a^{2})\,. 
\label{XXIX}
\end{equation}
The field strength also satifies the Bianchi identity 
\begin{equation}
D_\nu\,\widetilde G^{\mu \nu a} = 0\,.
\label{bianchi}
\end{equation}

Let us find the vacuum configurations in this theory. Using the
notations $G_a^{0i}=-{\cal E}_a^i$ and $G_a^{ij} = -\epsilon^{ij}_
{\hphantom{ij}k}{\cal B}_a^k$, the energy density is written as
\begin{equation}
\theta_{00}=\frac{1}{2}\left(({\cal E}_a^i)^2+({\cal B}_a^i)^2 +
  (D^0\phi_a)^2 + ( D^i \phi_a)^2 \right) + V (\phi)\,.
\label{XXX}
\end{equation} 
Note that $\theta_{00} \geq 0$, and it vanishes only if 
\begin{equation}
G_{a}^{\mu \nu} = 0, \quad D_{\mu} \phi = 0, \quad V(\phi) = 0\,.
\label{XXXI}
\end{equation} 
The first equation implies that in the vacuum, $W^a_\mu$ is pure
gauge and the last two equations define the Higgs vacuum. The
structure of the space of vacua is determined by $V(\phi)=0$ which
solves to $\phi^a= \phi^a_{vac}$ such that $|\phi_{vac}|=a$. The space
of Higgs vacua is therefore a two-sphere ($S^2$) of radius $a$ in the
field space. To formulate a perturbation theory, we have to choose one
of these vacua and hence, break the gauge symmetry spontaneously (this
is the usual Higgs mechanism). The part of the symmetry which keeps
this vacuum invariant, still survives and the corresponding unbroken 
generator is $\phi^c_{vac} T^c/a$. The gauge boson associated with
this generator is $A_\mu=\phi^c_{vac}W^c_\mu/a$ and the electric
charge operator for this surviving $U(1)$ is given by  
\begin{equation} 
Q = \hbar e \frac{\phi^c_{vac}T^c}{a}\,.
\label{XXXIII}
\end{equation}
If the group is compact, this charge is quantized. The perturbative
spectrum of the theory can be found by expanding $\phi^a$ around the
chosen vacuum as 
$$
\phi^{a} = \phi^{a}_{vac} + \phi '^{a}\,.
$$
A convenient choice is $\phi^c_{vac}=\delta^{c3}a$. The perturbative 
spectrum (which becomes manifest after choosing an appropriate gauge)
consists of a massive Higgs ($H$), a massless photon ($\gamma$) and
two charged massive bosons ($W^\pm$):  
$$
\begin{tabular}{|c|c|c|c|}
\hline
& Mass & Spin & Charge \\
\hline
$H$ & $a(2\lambda)^{\frac{1}{2}} \hbar$ & $0$ &$0$ \\
\hline
$\gamma$ & $0$ & $\hbar$ & $0$ \\
\hline
$W^{\pm}$ & $ae\hbar = aq$ & $\hbar$ & $\pm q = \pm e \hbar$ \\
\hline
\end{tabular}
$$
In the next section, we investigate the existence of monopoles
(non-perturbative states) in the Georgi-Glashow model.
